\documentclass[addpoints]{exam}
% Shared exam preamble (used with \documentclass[addpoints]{exam} in each test file)
\usepackage[margin=0.5in]{geometry}
\usepackage{tikz}
\usepackage{titling}
\usepackage{calc}
\usepackage{amsmath}
\usepackage{xparse}
\usepackage{multirow}
\usepackage{listings}
\usepackage{ifthen}
\usepackage{graphicx}
\usepackage{diagbox}
\usepackage{xcolor}
\usepackage{textcomp}
\usepackage{multicol}    % <<--- needed if you use multicols in the document
\usepackage{enumitem}

\newcommand{\overbar}[1]{\overline{#1}}

% SystemVerilog listings
\lstdefinelanguage{SystemVerilog}[]{Verilog}{
  morekeywords={
    logic,byte,int,shortint,longint,bit,enum,typedef,struct,union,
    always_comb,always_ff,always_latch,unique,priority,inside,with,
    foreach,rand,randc,assert,property,sequence,disable,fork,join_none,
    package,import,export,virtual,interface,modport
  }
}
\lstset{
  language=SystemVerilog,
  basicstyle=\ttfamily\small,
  keywordstyle=\color{blue},
  commentstyle=\color{gray},
  stringstyle=\color{red},
  numbers=left,
  numberstyle=\tiny,
  stepnumber=1,
  numbersep=5pt,
  showstringspaces=false,
  tabsize=2,
  breaklines=true,
  columns=fullflexible,
  literate={//§}{{\underline{\hspace{2cm}}}}3
}


\newcounter{subp}
\setcounter{subp}{0}

% 4-variable K-map for AB\CD (no parameters to avoid # in restricted mode)
\newcommand{\kmapABCD}{%
  {\renewcommand{\arraystretch}{1.1}\small
  \begin{tabular}{|c|c|c|c|c|}\hline
  AB\textbackslash CD & 00 & 01 & 11 & 10 \\ \hline
  00 &  &  &  &  \\ \hline
  01 &  &  &  &  \\ \hline
  11 &  &  &  &  \\ \hline
  10 &  &  &  &  \\ \hline
  \end{tabular}}%
}

% 2-input truth table: headers and 4 rows for 00,01,10,11
\newcommand{\truthtable}[3]{%
  {\renewcommand{\arraystretch}{1.1}\small
  \begin{tabular}{|c|c|c|}\hline
  #1 & #2 & #3 \\ \hline
  0  & 0  &    \\ \hline
  0  & 1  &    \\ \hline
  1  & 0  &    \\ \hline
  1  & 1  &    \\ \hline
  \end{tabular}}%
}

\newcommand{\asciitable}{%
  \begingroup
  \fontsize{8}{9}\selectfont
  \centering
  \begin{tabular}{*{8}{cc}}
  00 & NUL   & 01 & SOH   & 02 & STX   & 03 & ETX   \\
  04 & EOT   & 05 & ENQ   & 06 & ACK   & 07 & BEL   \\
  08 & BS    & 09 &       & 0A & LF    & 0B & VT    \\
  0C &       & 0D & CR    & 0E &       & 0F &      \\
  10 & DLE   & 11 & DC1   & 12 & DC2   & 13 & DC3   \\
  14 & DC4   & 15 & NAK   & 16 & SYN   & 17 & ETB   \\
  18 & CAN   & 19 & EM   & 1A & SUB   & 1B & ESC   \\
  1C & FS    & 1D & GS    & 1E & RS    & 1F & US    \\
  20 & space & 21 & !    & 22 & "    & 23 & \#   \\
  24 & \$    & 25 & \%   & 26 & \&   & 27 & '    \\
  28 & (     & 29 & )    & 2A & *    & 2B & +    \\
  2C & ,     & 2D & -    & 2E & .    & 2F & /    \\
  30 & 0     & 31 & 1    & 32 & 2    & 33 & 3    \\
  34 & 4     & 35 & 5    & 36 & 6    & 37 & 7    \\
  38 & 8     & 39 & 9    & 3A & :    & 3B & ;    \\
  3C & <     & 3D & =    & 3E & >    & 3F & ?    \\
  40 & @     & 41 & A    & 42 & B    & 43 & C    \\
  44 & D     & 45 & E    & 46 & F    & 47 & G    \\
  48 & H     & 49 & I    & 4A & J    & 4B & K    \\
  4C & L     & 4D & M    & 4E & N    & 4F & O    \\
  50 & P     & 51 & Q    & 52 & R    & 53 & S    \\
  54 & T     & 55 & U    & 56 & V    & 57 & W    \\
  58 & X     & 59 & Y    & 5A & Z    & 5B & [    \\
  5C & \textbackslash & 5D & ] & 5E & \textasciicircum & 5F & \_ \\
  60 & \textasciigrave & 61 & a & 62 & b & 63 & c \\
  64 & d     & 65 & e    & 66 & f    & 67 & g    \\
  68 & h     & 69 & i    & 6A & j    & 6B & k    \\
  6C & l     & 6D & m    & 6E & n    & 6F & o    \\
  70 & p     & 71 & q    & 72 & r    & 73 & s    \\
  74 & t     & 75 & u    & 76 & v    & 77 & w    \\
  78 & x     & 79 & y    & 7A & z    & 7B & \{   \\
  7C & |     & 7D & \}   & 7E & \textasciitilde & 7F &      \\
  \end{tabular}
  \endgroup
}

\newcommand{\test}[2]{\input{../#1.tex}\newcommand{\testname}{#2}}
\newcommand{\ul}[1]{\underline{\hspace{#1cm}}}
% Test title page - includes centered title, name/NetID fields, and newpage
% Optional logo can be set with \newcommand{\courselogo}{path/to/logo.png}
\newcommand{\maketesttitle}{%
  \begin{center}
  \ifdefined\courselogo
    \includegraphics[height=1.5cm]{\courselogo} \\[1em]
  \fi
  {\Huge \textbf{\coursenum{} \\[0.5em]
  \semester{} \\[0.5em]
  \coursename{} \\[1em]
  \testname}}
  \end{center}
  \vspace{50pt}
  \noindent Name: \ul{5} \\[50pt]
  \noindent NetID: \ul{2} \\[50pt]
  \newpage
}

\graphicspath{{../img/}{../figs/}}
\newcommand{\cpp}[1]{\lstinputlisting[language=C++]{code/#1}}
\newcommand{\cuda}[1]{\lstinputlisting[language=CUDA]{code/#1}}
\newcommand{\ccode}[1]{\lstinputlisting[language=C]{code/#1}}
\newcommand{\bash}[1]{\lstinputlisting[language=bash]{code/#1}}
\newcommand{\verilog}[1]{\lstinputlisting[language=SystemVerilog]{code/#1}}
\newcommand{\python}[1]{\lstinputlisting[language=Python]{code/#1}}


\title{ECE456 Network Centric Programming -- System Calls and Threads}
\date{}
\author{}

\begin{document}

\maketitle

\begin{questions}

	%------------------------------------------------
	\question[6] What is the output?

	\begin{lstlisting}[language=C]
#include <fcntl.h>
#include <unistd.h>

int main() {

    int fd = open("test.txt", O_WRONLY | O_CREAT, 0644);

    write(fd, "ABC", 3);
    write(fd, "DEF", 3);

    close(fd);
}
\end{lstlisting}

	Assume the file did not previously exist.

	What are the contents of \texttt{test.txt} after the program runs?

	\vspace{2cm}

	%------------------------------------------------
	\question[6]

	Fill in the blank so the program reads 4 bytes from a file.

	\begin{lstlisting}[language=C]
#include <fcntl.h>
#include <unistd.h>

int main() {

    int fd = open("data.bin", O_RDONLY);

    char buf[4];

    ________(fd, buf, 4);

    close(fd);
}
\end{lstlisting}

	\vspace{1cm}

	%------------------------------------------------
	\question[6]

	What error occurs? (single word)

	\begin{lstlisting}[language=C]
#include <fcntl.h>
#include <unistd.h>

int main() {

    int fd = open("file.txt", O_RDONLY);

    close(fd);

    char buf[10];
    read(fd, buf, 10);
}
\end{lstlisting}

	Error:

	\vspace{1cm}

	%------------------------------------------------
	\question[6]

	What system call creates a directory?

	\begin{lstlisting}[language=C]
#include <sys/stat.h>

int main() {

    ________("newdir", 0755);
}
\end{lstlisting}

	\vspace{1cm}

	%------------------------------------------------
	\question[6]

	What does the following program change?

	\begin{lstlisting}[language=C]
#include <unistd.h>

int main() {

    chdir("logs");
}
\end{lstlisting}

	Short answer:

	\vspace{2cm}

	%------------------------------------------------
	\question[10]

	Consider the following program.

	\begin{lstlisting}[language=C]
#include <unistd.h>
#include <stdio.h>

int main() {

    printf("A\n");

    fork();

    printf("B\n");
}
\end{lstlisting}

	List two possible outputs.

	\vspace{3cm}

	%------------------------------------------------
	\question[10]

	Consider:

	\begin{lstlisting}[language=C]
#include <unistd.h>
#include <stdio.h>

int main() {

    printf("1\n");

    pid_t pid = fork();

    if(pid == 0)
        printf("child\n");
    else
        printf("parent\n");
}
\end{lstlisting}

	How many processes run the program after the fork?

	Answer:

	\vspace{1cm}

	List two possible outputs.

	\vspace{3cm}

	%------------------------------------------------
	\question[12]

	Suppose a program writes:

	\begin{lstlisting}[language=C]
fprintf(stdout,"hello\n");
fprintf(stderr,"error\n");
\end{lstlisting}

	The program is executed as:

	\begin{verbatim}
./prog > out.txt 2> err.txt
\end{verbatim}

	Where does each line go?

	\begin{center}
		\begin{tabular}{l l}
			hello & \_\_\_\_\_\_\_\_\_\_\_ \\
			error & \_\_\_\_\_\_\_\_\_\_\_
		\end{tabular}
	\end{center}

	%------------------------------------------------
	\question[12]

	Explain what the following command does.

	\begin{verbatim}
./prog > output.txt 2>&1
\end{verbatim}

	Short answer:

	\vspace{3cm}

	%------------------------------------------------
	\question[15]

	Write code that:

	\begin{itemize}
		\item opens a file called \texttt{numbers.txt}
		\item writes the numbers 0--9 (one per line)
		\item closes the file
	\end{itemize}

	Use only:

	\begin{itemize}
		\item \texttt{open}
		\item \texttt{write}
		\item \texttt{close}
	\end{itemize}

	\vspace{7cm}

\end{questions}

\end{document}